% chapters/06/checklist.tex
\bagian{Checklist Pencapaian Kompetensi}

Bab ini menuntaskan \SUBCPMK{1.3}. Centang setiap butir berikut hanya apabila
Anda benar-benar telah melakukannya atau dapat menjelaskannya.

\begin{checklistbox}
\begin{checklist}
  \item {\bfseries (Indikator a)} Saya dapat menulis integrasi Euler dua tingkat,
        yaitu memperbarui kecepatan dengan percepatan lalu posisi dengan
        kecepatan.
  \item Saya dapat menjelaskan mengapa urutan itu tidak boleh dibalik.
  \item Saya dapat menjelaskan tanda gravitasi dan tanda kecepatan lompatan
        pada koordinat layar.
  \item Saya dapat menghentikan benda di lantai dan menjelaskan mengapa
        pemeriksaannya memakai tanda ketidaksamaan.
  \item Saya dapat menjelaskan mengapa posisi harus diperbaiki sebelum
        kecepatan diubah ketika benda menyentuh lantai.
  \item {\bfseries (Indikator b)} Saya dapat menulis lompatan sebagai impuls,
        yaitu penggantian kecepatan tegak dengan satu baris.
  \item Lompatan pada program saya hanya bekerja ketika benda menyentuh lantai,
        sehingga pemain tidak dapat melompat di udara.
  \item Saya dapat menjelaskan mengapa lompatan dibaca sebagai peristiwa, dan
        menyebutkan gejala yang muncul apabila dibaca sebagai keadaan tombol.
  \item Saya dapat menghitung kecepatan lompatan yang diperlukan untuk
        memperoleh tinggi lompatan tertentu.
  \item {\bfseries (Indikator c)} Saya dapat menulis pantulan pada lantai dan
        pada dinding dengan koefisien restitusi.
  \item Saya dapat menjelaskan mengapa tinggi pantulan menyusut sebesar kuadrat
        koefisien restitusi.
  \item Saya dapat menghentikan pantulan yang terlalu kecil agar benda tidak
        bergetar di permukaan lantai.
  \item Saya telah mengukur tinggi lompatan pada tiga laju bingkai dan
        membandingkannya dengan hitungan.
  \item Catatan praktikum saya memuat hasil percobaan koefisien restitusi pada
        bola.
\end{checklist}
\end{checklistbox}
